% ==============================================================================
% Plantilla Institucional de Trabajo Terminal — UPIIZ IPN
% Archivo: capitulos/protocolo/01-objetivos.tex
% ==============================================================================
% ESTRUCTURA NORMATIVA DE PROTOCOLO: CAPÍTULO 1 - OBJETIVOS DEL PROYECTO
% ==============================================================================
% Directrices institucionales (Anexo 1):
% - Iniciar la redacción de cada objetivo estrictamente con un verbo en infinitivo.
% - El Objetivo General debe definir con claridad qué se desarrollará, mediante qué
%   enfoque o tecnología y con qué propósito o finalidad.
% - Los Objetivos Específicos (o particulares) deben ser como MÁXIMO 4.
% - Los objetivos específicos representan metas concretas y secuenciales para
%   alcanzar el objetivo general, sin condicionar la naturaleza del proyecto.
% - Debe incluirse la firma de conformidad de alumnos y asesores.
% ==============================================================================

\chapter{Objetivos del proyecto}
\label{cap:proto-objetivos}

En este capítulo se formulan las metas del Trabajo Terminal propuesto, estableciendo el alcance funcional y los compromisos técnicos del proyecto.

% ------------------------------------------------------------------------------
\vspace{-0.2cm}
\section{Objetivo general}
\label{sec:proto-objetivo-general}

% [INSTRUCCIÓN PARA EL ALUMNO]:
% Redacte en un solo párrafo el Objetivo General de su propuesta.
% Inicie con un verbo en infinitivo que defina con precisión la meta global del proyecto:
% [Verbo en infinitivo] + [sistema, dispositivo, proceso o software a desarrollar] +
% [enfoque técnico, metodológico o disciplinar] + [propósito o solución planteada].

[Inserte aquí la redacción del Objetivo General de su propuesta de Trabajo Terminal].

% ------------------------------------------------------------------------------
\vspace{-0.2cm}
\section{Objetivos específicos o particulares}
\label{sec:proto-objetivos-especificos}

% [INSTRUCCIÓN PARA EL ALUMNO]:
% Formule un MÁXIMO DE 4 objetivos específicos (o particulares) ordenados cronológica
% y lógicamente. Cada objetivo debe redactarse en infinitivo y corresponder a una fase
% medible para el cumplimiento del objetivo general:

\vspace{-0.2cm}
\begin{enumerate}\setlength{\itemsep}{0pt}\setlength{\parskip}{0pt}\setlength{\parsep}{0pt}
    \item \textbf{Objetivo específico 1:} [Verbo en infinitivo] los requerimientos, especificaciones técnicas y análisis inicial del sistema propuesto.
    \item \textbf{Objetivo específico 2:} [Verbo en infinitivo] el diseño, modelado y simulación de los componentes o subsistemas requeridos.
    \item \textbf{Objetivo específico 3:} [Verbo en infinitivo] la implementación, desarrollo o ensamble de los módulos constitutivos del proyecto.
    \item \textbf{Objetivo específico 4:} [Verbo en infinitivo] la integración final, pruebas y evaluación experimental del desempeño del sistema.
\end{enumerate}

% ------------------------------------------------------------------------------
% CONFORMIDAD DE OBJETIVOS (FIRMAS DE ALUMNOS Y ASESORES)
% ------------------------------------------------------------------------------
\vspace{0.1cm}
\noindent
\begin{minipage}{\textwidth}
    \centering
    {\fontsize{10.5}{12}\bfseries\sffamily Conformidad de Objetivos del Proyecto}\\[0.15em]
    {\fontsize{8.5}{10}\sffamily Con la firma de esta sección, los alumnos y asesores avalan formalmente los objetivos aquí establecidos.}\\[0.3em]
    \bloqueFirmasProtocolo[0.2cm]
\end{minipage}
